% The classical construction: Besselian elements (one frame; input file).
% Input by math_slides.tex, unless the main document shows it earlier and sets \BesselianShown
% (the talk: at the start of "How it works"). Found via \input@path from other folders.

% the classical construction (Wikimedia Commons figure, CC BY-SA 3.0; credit as an unnumbered footnote)
\begin{frame}{The classical construction: Besselian elements}
  \begin{columns}[c]
    \begin{column}{0.48\textwidth}
      \centering
      \includegraphics[height=0.74\textheight]{BesselianElementsForTotalEclipse3D-en.pdf}
    \end{column}
    \begin{column}{0.50\textwidth}
      \small
      \begin{itemize}\setlength{\itemsep}{1ex}
        \item Solar eclipse: the Moon's shadow is a \alert{cone} (penumbra, umbra); the fundamental plane
              ($x$, $y$) through the Earth's centre, perpendicular to the shadow axis
        \item Occultation of a star: the star is at infinity, so the asteroid's shadow is a \alert{cylinder}
              of the asteroid's size
        \item Our solver: the plane through \alert{the observer}, the asteroid's light time from the site
              (slide FP)
      \end{itemize}
    \end{column}
  \end{columns}
  {\renewcommand{\thefootnote}{}\footnotetext{\tiny Figure: By G.~Pieper, enhanced by User:Iridos, Uploader: Cactus26
   - Own work, CC BY-SA 3.0, \url{https://commons.wikimedia.org/w/index.php?curid=12020673}}}
\end{frame}
